\begin{abstract}
Traditional music recommendation systems primarily rely on collaborative and content-based filtering, which often struggle with the "filter bubble" effect and lack explainability. With the advent of Large Language Models (LLMs), there is a growing interest in leveraging their semantic reasoning capabilities for personalized recommendation. In this work, we examine whether LLMs can outperform established recommendation algorithms and human-curated benchmarks when provided with a user's listening history. We simulate a "Spotify Data Export" workflow using the \musicrs benchmark and real-world listening sessions, evaluating recommendations from \claudesonnet and \gptfour. Our results show that LLMs matched or outperformed human-curated recommendations in 80\% of test cases, particularly excelling in cross-genre reasoning and thematic coherence. We demonstrate that while traditional models maintain an edge in processing high-velocity behavioral data, LLMs offer a paradigm shift towards context-aware, explainable music discovery.
\end{abstract}
